% ============================================================================
% TCC — Predição de Direção de Preços de Ações com NLP e Sentimento de Notícias
% Formato: abnTeX2 (ABNT NBR 14724:2011)
% Compilar: pdflatex tcc.tex && bibtex tcc && pdflatex tcc.tex && pdflatex tcc.tex
% ============================================================================
\documentclass[
    12pt,
    openright,
    oneside,
    a4paper,
    english,
    brazil
]{abntex2}

% --- Pacotes ---
\usepackage[utf8]{inputenc}
\usepackage[T1]{fontenc}
\usepackage{lmodern}
\usepackage{amsmath,amssymb}
\usepackage{graphicx}
\usepackage{booktabs}
\usepackage{tabularx}
\usepackage{multirow}
\usepackage{float}
\usepackage{longtable}
\usepackage{microtype}
\usepackage{hyperref}
\usepackage[brazilian,hyperpageref]{backref}
\usepackage[alf]{abntex2cite}
\usepackage{xcolor}
\usepackage{listings}
\usepackage{enumitem}

% --- Configurações de listings ---
\lstset{
    basicstyle=\ttfamily\small,
    breaklines=true,
    frame=single,
    numbers=left,
    numberstyle=\tiny\color{gray},
    keywordstyle=\color{blue},
    commentstyle=\color{green!50!black},
    stringstyle=\color{red!60!black},
}

% --- Configurações de backref ---
\renewcommand{\backrefpagesname}{Citado na(s) página(s):~}
\renewcommand{\backref}{}
\renewcommand*{\backrefalt}[4]{
    \ifcase #1
        Nenhuma citação no texto.
    \or
        Citado na página #2.
    \else
        Citado #1 vezes nas páginas #2.
    \fi
}

% --- Informações do documento ---
\titulo{Predição de Direção de Preços de Ações Brasileiras\\com Sentimento de Notícias Financeiras:\\Pipeline, Ilusão e Autocorreção Metodológica}
\autor{André Takeo Loschner Fujiwara}
\local{São Paulo}
\data{2026}
\orientador{Prof.\ Eric Bacconi Gonçalves}
\instituicao{%
    PUC-SP
    \par Curso de Ciência de Dados e Inteligência Artificial
}
\tipotrabalho{Trabalho de Conclusão de Curso}
\preambulo{Trabalho de Conclusão de Curso apresentado ao Curso de Ciência de Dados e Inteligência Artificial como requisito parcial para a obtenção do grau de Bacharel.}

% --- Compila o índice ---
\makeindex

% ============================================================================
\begin{document}

\selectlanguage{brazil}
\frenchspacing

% --- Capa ---
\imprimircapa

% --- Folha de rosto ---
\imprimirfolhaderosto*

% ============================================================================
% ELEMENTOS PRÉ-TEXTUAIS
% ============================================================================

% --- Resumo ---
\begin{resumo}
Este trabalho investiga se notícias financeiras brasileiras podem melhorar a predição de direção de preços de ações da B3, utilizando como estudo de caso três ações de grande capitalização: ITUB4 (Itaú Unibanco), PETR4 (Petrobras) e VALE3 (Vale). O \textit{pipeline} proposto coleta artigos do portal InfoMoney, extrai representações de sentimento via FinBERT-PT-BR --- um modelo BERT especializado em textos financeiros em português --- e combina as \textit{features} resultantes com dados de mercado do Yahoo Finance para treinar classificadores binários (sobe/desce).

A dissertação documenta uma autocorreção metodológica incomum: sob avaliação por janela única (\textit{walk-forward split} 70/15/15), o modelo Transformer atinge ROC-AUC de 0,709, aparentemente demonstrando que o sentimento financeiro específico supera \textit{embeddings} genéricos de alta dimensionalidade. Porém, uma investigação sistemática com 1.435 execuções de modelo --- incluindo análise \textit{multi-seed} (20 sementes), validação cruzada \textit{expanding-window} (até 52 \textit{folds}), testes de Wilcoxon \textit{signed-rank}, intervalos de confiança \textit{bootstrap} e \textit{ablation} de \textit{features} --- revela que o resultado é um artefato: o Transformer exibe colapso bimodal (AUC oscilando entre 0,08 e 0,93 conforme a semente), e sob avaliação multi-fold o sentimento adiciona $\Delta = +0{,}003$ ao \textit{baseline} autoregressivo (Wilcoxon $p = 0{,}49$, IC 95\% $[-0{,}012;\ +0{,}018]$).

A contribuição central é a demonstração empírica, em dados brasileiros, de como a avaliação por janela única pode inverter a conclusão qualitativa de um estudo em \textit{machine learning} financeiro --- transformando um achado aparentemente positivo em um resultado nulo --- e a proposição de protocolos mínimos de avaliação para mitigar esse viés.

\textbf{Palavras-chave:} Predição de preços. Sentimento financeiro. FinBERT. Avaliação metodológica. Séries temporais financeiras. B3.
\end{resumo}

% --- Abstract ---
\begin{resumo}[Abstract]
\begin{otherlanguage*}{english}
This work investigates whether Brazilian financial news can improve stock price direction prediction for B3 equities, using three large-cap stocks as case studies: ITUB4 (Itaú Unibanco), PETR4 (Petrobras), and VALE3 (Vale). The proposed pipeline collects articles from the InfoMoney portal, extracts sentiment representations via FinBERT-PT-BR --- a BERT model fine-tuned on Brazilian Portuguese financial text --- and combines the resulting features with Yahoo Finance market data to train binary classifiers (up/down).

The dissertation documents an uncommon methodological self-correction: under single-window evaluation (70/15/15 walk-forward split), the Transformer model achieves ROC-AUC of 0.709, apparently demonstrating that domain-specific financial sentiment outperforms high-dimensional generic embeddings. However, a systematic investigation comprising 1,435 model runs --- including multi-seed analysis (20 seeds), expanding-window cross-validation (up to 52 folds), Wilcoxon signed-rank tests, bootstrap confidence intervals, and feature ablation --- reveals the result to be an artifact: the Transformer exhibits bimodal collapse (AUC ranging from 0.08 to 0.93 depending on seed), and under multi-fold evaluation sentiment adds $\Delta = +0.003$ over the autoregressive baseline (Wilcoxon $p = 0.49$, 95\% CI $[-0.012,\ +0.018]$).

The central contribution is an empirical demonstration, on Brazilian data, of how single-window evaluation can reverse the qualitative conclusion of a financial machine learning study --- turning an apparently positive finding into a null result --- along with minimal evaluation protocols to mitigate this bias.

\textbf{Keywords:} Price prediction. Financial sentiment. FinBERT. Methodological evaluation. Financial time series. B3.
\end{otherlanguage*}
\end{resumo}

% --- Lista de figuras ---
\pdfbookmark[0]{\listfigurename}{lof}
\listoffigures*
\cleardoublepage

% --- Lista de tabelas ---
\pdfbookmark[0]{\listtablename}{lot}
\listoftables*
\cleardoublepage

% --- Lista de abreviaturas ---
\begin{siglas}
    \item[AUC] \textit{Area Under the ROC Curve}
    \item[B3] Brasil, Bolsa, Balcão
    \item[BiLSTM] \textit{Bidirectional Long Short-Term Memory}
    \item[CI] Intervalo de Confiança (\textit{Confidence Interval})
    \item[CV] Validação Cruzada (\textit{Cross-Validation})
    \item[ECE] \textit{Expected Calibration Error}
    \item[FinBERT] \textit{Financial BERT}
    \item[MDE] Tamanho Mínimo de Efeito Detectável
    \item[NLP] Processamento de Linguagem Natural
    \item[OHLCV] \textit{Open, High, Low, Close, Volume}
    \item[PCA] Análise de Componentes Principais
    \item[ROC] \textit{Receiver Operating Characteristic}
    \item[SHAP] \textit{SHapley Additive exPlanations}
    \item[TCN] \textit{Temporal Convolutional Network}
    \item[XGB] XGBoost (\textit{eXtreme Gradient Boosting})
\end{siglas}

% --- Sumário ---
\pdfbookmark[0]{\contentsname}{toc}
\tableofcontents*
\cleardoublepage

% ============================================================================
% ELEMENTOS TEXTUAIS
% ============================================================================
\textual

% ============================================================================
\chapter{Introdução}\label{cap:introducao}
% ============================================================================

A predição de direção de preços de ações é um problema central em finanças quantitativas, situado na interseção entre a hipótese de eficiência de mercado \cite{fama1970} e a evidência empírica de previsibilidade parcial em horizontes curtos \cite{lo2004}. Com o avanço das técnicas de processamento de linguagem natural (NLP), uma linha de pesquisa crescente investiga se informações textuais --- notícias, relatórios de analistas, publicações em redes sociais --- podem ser incorporadas como \textit{features} preditivas em modelos de \textit{machine learning} para complementar dados tradicionais de preço e volume \cite{xu2018,araci2019}.

No contexto brasileiro, a pesquisa nesta área é ainda incipiente. A B3 (Brasil, Bolsa, Balcão) apresenta características que a distinguem de mercados desenvolvidos: menor liquidez em ações fora do índice principal, concentração de informação em poucas fontes jornalísticas especializadas, e um ecossistema de NLP em português que só recentemente passou a dispor de modelos de linguagem específicos para o domínio financeiro, como o FinBERT-PT-BR.

\section{Problema de pesquisa}

O problema central pode ser formulado como: \textbf{notícias financeiras brasileiras melhoram a predição de direção de preço (sobe/desce) de ações da B3?} Esta questão desdobra-se em duas dimensões:

\begin{enumerate}[label=(\alph*)]
    \item \textbf{Representação textual:} qual forma de codificar o conteúdo de notícias --- \textit{embeddings} genéricos de alta dimensionalidade ou sentimento financeiro específico de baixa dimensionalidade --- produz \textit{features} mais informativas para a predição?
    \item \textbf{Robustez metodológica:} os resultados obtidos sob protocolos de avaliação padrão da literatura (\textit{walk-forward split} único) sobrevivem a escrutínio com validação cruzada \textit{expanding-window}, múltiplas sementes e testes estatísticos formais?
\end{enumerate}

\section{Objetivos}

\subsection{Objetivo geral}

Investigar empiricamente se \textit{features} de sentimento extraídas de notícias financeiras brasileiras adicionam poder preditivo a modelos de classificação binária de direção de preço de ações da B3, sob avaliação metodologicamente rigorosa.

\subsection{Objetivos específicos}

\begin{enumerate}
    \item Construir um \textit{pipeline} de coleta, processamento e integração de notícias financeiras do portal InfoMoney com dados de mercado do Yahoo Finance.
    \item Comparar duas estratégias de representação textual: \textit{embeddings} genéricos de 1.024 dimensões (Ollama \texttt{qwen3-embedding:4b}) e sentimento financeiro de 5 dimensões (FinBERT-PT-BR).
    \item Treinar e avaliar modelos de classificação binária (BiLSTM, Transformer, XGBoost, TCN) sobre as representações textuais combinadas com \textit{features} de preço.
    \item Aplicar protocolos de avaliação rigorosa --- \textit{bootstrap} CI, análise \textit{multi-seed}, validação cruzada \textit{expanding-window}, \textit{ablation} de \textit{features} --- para determinar a robustez dos resultados.
    \item Documentar e analisar as discrepâncias entre avaliação por janela única e avaliação multi-fold, contribuindo para a literatura de avaliação metodológica em ML financeiro.
\end{enumerate}

\section{Justificativa}

A relevância deste trabalho reside em três aspectos. Primeiro, a aplicação de NLP financeiro ao mercado brasileiro é pouco explorada, e a disponibilidade recente de modelos como o FinBERT-PT-BR abre oportunidades de investigação em português. Segundo, a literatura de \textit{deep learning} aplicado a finanças sofre de um problema metodológico amplamente documentado --- a avaliação por janela única em séries não-estacionárias \cite{bailey2014,lopezdeprado2018} --- mas raramente demonstrado empiricamente com o grau de detalhe que este trabalho propõe. Terceiro, a autocorreção metodológica documentada nos Capítulos~\ref{cap:sentimento} e~\ref{cap:investigacao} constitui, por si só, uma contribuição pedagógica sobre as armadilhas da avaliação de modelos em finanças.

\section{Estrutura do trabalho}

O trabalho está organizado em seis capítulos. O Capítulo~\ref{cap:pipeline} descreve o \textit{pipeline} completo: coleta de notícias, engenharia de \textit{features} e treinamento inicial com \textit{embeddings} genéricos. O Capítulo~\ref{cap:sentimento} introduz o FinBERT-PT-BR e documenta o resultado inicialmente promissor (AUC = 0,709). O Capítulo~\ref{cap:investigacao} conduz a investigação metodológica que revisa substancialmente esse resultado. O Capítulo~\ref{cap:conclusao} sintetiza as conclusões e propõe direções futuras.


% ============================================================================
\chapter{Pipeline: Coleta, Engenharia de Features e Modelos Iniciais}\label{cap:pipeline}
% ============================================================================

Este capítulo descreve as três primeiras etapas do \textit{pipeline}: coleta de notícias financeiras, engenharia de \textit{features} (mercado e texto) e treinamento inicial de modelos com \textit{embeddings} genéricos.

\section{Coleta de notícias financeiras}\label{sec:coleta}

O corpus textual foi construído a partir do portal InfoMoney, uma das principais fontes de notícias financeiras no Brasil. O módulo \texttt{extractor.py} acessa a REST API do WordPress do portal (\texttt{/wp-json/wp/v2/posts}), realizando buscas paginadas (100 artigos por página) para cada \textit{ticker} analisado.

Cada requisição HTTP possui mecanismo de retentativa: até 3 tentativas com \textit{backoff} linear (1\,s, 2\,s), tratando erros de \textit{timeout}, conexão e códigos de status 429/5xx. Os artigos passam por um \textit{pipeline} de pré-processamento: remoção de \textit{tags} HTML, decodificação de entidades, normalização de \textit{whitespace} e extração de \textit{keywords} do \textit{schema} Yoast SEO (JSON-LD). A extração de múltiplos \textit{tickers} ocorre em paralelo via \texttt{ThreadPoolExecutor}.

\begin{table}[htb]
\centering
\caption{Volume de artigos coletados por ativo.}\label{tab:artigos}
\begin{tabular}{lrr}
\toprule
\textbf{Ativo} & \textbf{Artigos} & \textbf{Período} \\
\midrule
ITUB4 (Itaú Unibanco) & 2.572 & 2009--2026 \\
PETR4 (Petrobras)     & 1.775 & 2009--2026 \\
VALE3 (Vale)          & 1.525 & 2009--2026 \\
\midrule
\textbf{Total}        & \textbf{5.872} & \\
\bottomrule
\end{tabular}
\end{table}

\section{Engenharia de features}\label{sec:features}

\subsection{Features de mercado}

O módulo \texttt{yahoo\_finance.py} utiliza a biblioteca \texttt{yfinance} para obter dados OHLCV. A partir dos dados brutos, são calculadas 11 \textit{features} técnicas por dia:

\begin{itemize}
    \item Preço de fechamento (\texttt{Close}) e volume negociado (\texttt{Volume});
    \item Retorno diário percentual (\texttt{return});
    \item Médias móveis de 7 e 21 dias (\texttt{ma7}, \texttt{ma21});
    \item Desvio-padrão de 21 dias como \textit{proxy} de volatilidade (\texttt{std21});
    \item Valores defasados do fechamento (\texttt{lag\_1} a \texttt{lag\_5}).
\end{itemize}

\subsection{Embeddings de notícias}

O módulo \texttt{news\_embedder.py} transforma cada artigo em um vetor de 1.024 dimensões usando o modelo \texttt{qwen3-embedding:4b} via Ollama. A entrada é a concatenação de título, resumo e conteúdo (truncado em 50.000 caracteres). Quando há múltiplas notícias no mesmo dia, os \textit{embeddings} são agregados via média ponderada por recência: notícias mais recentes recebem peso maior (pesos lineares 1, 2, 3\ldots).

\subsection{Merge temporal}

As \textit{features} de preço e os \textit{embeddings} diários são combinados via \textit{left join} por data. Para dias de mercado sem notícias publicadas, aplica-se \textit{forward-fill} nos \textit{embeddings}. O \textit{dataset} resultante contém 1.227 dias $\times$ 1.035 \textit{features} (11 de preço + 1.024 de \textit{embeddings}), cobrindo o período de abril de 2021 a março de 2026 para o ITUB4.

\section{Modelos iniciais com embeddings genéricos}\label{sec:modelos_iniciais}

O alvo binário é definido como 1 se $\text{Close}[t+21] > \text{Close}[t]$, resultando em desbalanceamento de 59\% ``Sobe'' / 41\% ``Desce''. Os 1.024 \textit{embeddings} são reduzidos para 32 componentes via PCA (61,4\% de variância explicada). A validação segue \textit{split} cronológico 70/15/15 sem \textit{shuffle}. Quatro modelos foram avaliados:

\begin{table}[htb]
\centering
\caption{Resultados dos modelos treinados com \textit{embeddings} genéricos Ollama (Etapa~3).}\label{tab:etapa3}
\begin{tabular}{lcc}
\toprule
\textbf{Modelo} & \textbf{Arquitetura} & \textbf{ROC-AUC} \\
\midrule
BiLSTM Original  & 2 camadas, 128 \textit{hidden}, 30\% \textit{dropout} & 0,443 \\
BiLSTM Reduzido  & 1 camada, 32 \textit{hidden}, 50\% \textit{dropout}  & 0,505 \\
Transformer      & 2 camadas, 4 cabeças, $d_\text{model}=64$              & 0,568 \\
\textbf{XGBoost} & 300 árvores, profundidade 4                            & \textbf{0,610} \\
\bottomrule
\end{tabular}
\end{table}

Todos os modelos apresentaram desempenho próximo do acaso, indicando que a combinação de \textit{features} de preço com \textit{embeddings} genéricos de alta dimensionalidade não é suficiente para predição de direção. A interpretação inicial atribuiu o resultado a duas hipóteses: (A) eficiência de mercado ou (B) inadequação da representação textual. O Capítulo~\ref{cap:sentimento} testa a hipótese~(B).


% ============================================================================
\chapter{Sentimento Financeiro Específico via FinBERT-PT-BR}\label{cap:sentimento}
% ============================================================================

\section{Hipótese}

Os resultados do Capítulo~\ref{cap:pipeline} sugeriram que os \textit{embeddings} genéricos contêm ruído semântico irrelevante para o domínio financeiro. Este capítulo testa se uma representação compacta e específica --- 5 \textit{features} de sentimento financeiro extraídas com o FinBERT-PT-BR --- pode ser mais informativa.

\section{O modelo FinBERT-PT-BR}

O FinBERT-PT-BR é um modelo BERT pré-treinado e \textit{fine-tuned} em corpora financeiros brasileiros. Para cada artigo, produz três \textit{logits} correspondentes às classes POSITIVO, NEGATIVO e NEUTRO. Diferentemente de modelos de sentimento genéricos, o FinBERT-PT-BR foi calibrado para vocabulário financeiro: termos como ``déficit'', ``rentabilidade'', ``\textit{guidance}'' e ``\textit{downgrade}'' recebem peso semântico adequado ao domínio.

\section{Pipeline de extração}

A extração de sentimento segue os passos: (1)~carregamento dos artigos coletados na Etapa~1; (2)~inferência em \textit{batch} (32 artigos) com entrada \texttt{título + resumo} truncada em 512 \textit{tokens}; (3)~persistência dos \textit{logits} brutos e classe predita por artigo; (4)~agregação diária em 5 \textit{features}: \texttt{n\_articles}, \texttt{mean\_logit\_pos}, \texttt{mean\_logit\_neg}, \texttt{mean\_logit\_neu} e \texttt{mean\_sentiment}; (5)~junção temporal com \textit{features} de preço via \textit{left join} com \textit{forward-fill}.

O \textit{dataset} resultante contém 1.207 dias $\times$ 16 \textit{features} (11 de preço + 5 de sentimento).

\textbf{Granularidade temporal e risco de \textit{look-ahead}.} Os artigos possuem \textit{timestamps} completos (ISO~8601), mas a agregação diária descarta o componente horário. Artigos publicados após o fechamento da B3 ($\sim$17h00 BRT) são atribuídos ao dia $t$, embora seu efeito possa se manifestar apenas em $t+1$. O deslocamento pode exceder um dia quando há dias consecutivos sem publicação. Três fatores atenuam o impacto: (1)~o horizonte de 5--21 dias dilui deslocamentos de 1 dia; (2)~a média diária dilui artigos individuais; (3)~o \textit{forward-fill} opera na direção causal correta.

\section{Resultados sob walk-forward split único}

\begin{table}[htb]
\centering
\caption{Comparação Etapa~3 (\textit{embeddings}) vs Etapa~4 (sentimento FinBERT).}\label{tab:etapa4}
\begin{tabular}{lccc}
\toprule
\textbf{Modelo} & \textbf{AUC Etapa~3} & \textbf{AUC Etapa~4} & $\boldsymbol{\Delta}$ \\
\midrule
BiLSTM Original  & 0,443 & 0,500 & +0,057 \\
BiLSTM Reduzido  & 0,505 & 0,477 & $-$0,028 \\
XGBoost          & 0,610 & 0,670 & +0,060 \\
\textbf{Transformer} & 0,568 & \textbf{0,709} & \textbf{+0,141} \\
\bottomrule
\end{tabular}
\end{table}

O Transformer atinge ROC-AUC = 0,709, o melhor resultado do trabalho sob este protocolo. Métricas adicionais: acurácia 76,3\% (177 amostras de teste), F1(Sobe) = 0,85, \textit{precision}(Desce) = 1,00 mas \textit{recall}(Desce) = 0,20 --- o modelo prevê ``Desce'' apenas 11 vezes em 177 amostras.

Três observações qualitativas motivam a investigação do Capítulo~\ref{cap:investigacao}: (1)~a matriz de confusão é altamente assimétrica; (2)~a acurácia é próxima da proporção da classe majoritária ($\sim$69\%); (3)~nem intervalos de confiança nem múltiplas sementes foram reportados.


% ============================================================================
\chapter{Investigação Metodológica e o Viés da Avaliação por Janela Única}\label{cap:investigacao}
% ============================================================================

\section{Motivação}

A vulnerabilidade da avaliação por janela única em séries financeiras é documentada na literatura: \citeonline{bailey2014} demonstraram formalmente o \textit{backtest overfitting}; \citeonline{lopezdeprado2018} propôs validação cruzada \textit{purged/embargo}; \citeonline{cawley2010} documentaram o viés de seleção de modelo. A adoção do \textit{walk-forward split} único nos capítulos anteriores seguiu o padrão predominante na literatura \cite{fischer2018,xu2018,araci2019}. Este capítulo demonstra empiricamente, com 1.435 execuções, o alcance das distorções que esse protocolo pode produzir.

\section{Protocolo geral}

Todos os experimentos compartilham o mesmo \textit{dataset} base, o mesmo alvo binário e a mesma arquitetura Transformer do Capítulo~\ref{cap:sentimento}. Os utilitários reutilizáveis (\textit{split walk-forward}, \textit{bootstrap} CI, geração de alvo) estão em \texttt{eval\_utils.py}. A Tabela~\ref{tab:experimentos} sumariza os 1.435 treinamentos realizados.

\begin{table}[htb]
\centering
\caption{Experimentos do Capítulo~\ref{cap:investigacao}: 1.435 treinamentos de modelo.}\label{tab:experimentos}
\begin{tabular}{llr}
\toprule
\textbf{Experimento} & \textbf{Notebook} & \textbf{Runs} \\
\midrule
Baseline autoregressivo          & \texttt{dumb\_baseline.ipynb}                & 4   \\
Baselines ingênuos               & \texttt{naive\_baselines.ipynb}              & 3   \\
Controle de dimensionalidade     & \texttt{dimensionality\_control.ipynb}       & 20  \\
Bootstrap CI (Etapa~4)           & \texttt{bootstrap\_stage4.ipynb}             & 1   \\
Multi-seed em ITUB4              & \texttt{multi\_seed.ipynb}                   & 40  \\
Multi-seed $\times$ multi-ticker & \texttt{multi\_seed\_multi\_ticker.ipynb}     & 120 \\
Ensemble + backtest              & \texttt{ensemble\_backtest.ipynb}             & 20  \\
Expanding-window CV              & \texttt{expanding\_window\_cv.ipynb}          & 145 \\
Varredura de horizontes          & \texttt{horizon\_sweep.ipynb}                & 6   \\
VALE3 \textit{deep-dive}        & \texttt{vale3\_deepdive.ipynb}               & 880 \\
Ablation PRICE/SENT/PRICE+SENT  & \texttt{ablation\_price\_vs\_sentiment.ipynb} & 225 \\
Poder estatístico                & \texttt{power\_analysis.ipynb}               & --- \\
Validação TCN Stage~5b           & \texttt{tcn\_validation.ipynb}               & 45  \\
\midrule
\textbf{Total}                   &                                              & \textbf{$\sim$1.510} \\
\bottomrule
\end{tabular}
\end{table}

\section{Baseline autoregressivo e baselines ingênuos}

Um XGBoost treinado em 5 \textit{features} autoregressivas de preço (\texttt{return}, \texttt{lag\_1}, \texttt{lag\_5}, \texttt{Volume}, \texttt{std21}), sem qualquer informação textual, atinge AUC = 0,658 [0,565; 0,744]. O \textit{gap} contra o resultado original (0,709) é de apenas 0,051 pontos.

Para contextualizar, três preditores verdadeiramente ingênuos foram avaliados (Tabela~\ref{tab:naive}).

\begin{table}[htb]
\centering
\caption{Hierarquia de \textit{baselines}: do ingênuo ao autoregressivo.}\label{tab:naive}
\begin{tabular}{llc}
\toprule
\textbf{Preditor} & \textbf{Descrição} & \textbf{AUC [IC 95\%]} \\
\midrule
Classe majoritária     & Sempre prevê ``Sobe''           & 0,500 [0,500; 0,500] \\
\textit{Coin flip}     & $P(\text{Sobe}) = \text{prior}$ & 0,500 [0,500; 0,500] \\
Persistência ($h=21$)  & Direção repete                  & 0,474 [0,405; 0,543] \\
\textbf{Autoregressivo (XGB)} & \textbf{5 features preço} & \textbf{0,658 [0,565; 0,744]} \\
\bottomrule
\end{tabular}
\end{table}

\subsection{Controle de dimensionalidade}

A comparação Etapa~3 (1.024~dim) $\to$ Etapa~4 (5~dim) confunde representação com dimensionalidade. XGBoost treinado em 20 subconjuntos aleatórios de 5 dimensões Ollama produz AUC médio de $0{,}509 \pm 0{,}057$ --- próximo do acaso e inferior ao FinBERT (0,670), confirmando que a melhoria reflete especificidade de domínio, não mera redução dimensional.

\section{Reprodutibilidade e variância de semente}

Sob recompilação com semente 42, o Transformer atinge AUC = 0,442 [0,358; 0,528] --- uma diferença de 0,267 pontos, equivalente a $22\times$ o desvio-padrão do \textit{baseline} (0,012). Sob 20 sementes distintas em ITUB4:

\begin{table}[htb]
\centering
\caption{Distribuição de AUC sobre 20 sementes (ITUB4, janela única).}\label{tab:multiseed}
\begin{tabular}{lcccccc}
\toprule
\textbf{Modelo} & \textbf{Média} & \textbf{Std} & \textbf{Mediana} & \textbf{Min} & \textbf{Max} \\
\midrule
Transformer + FinBERT       & 0,686 & \textbf{0,261} & 0,802 & 0,080 & 0,931 \\
Baseline autoregressivo     & 0,641 & 0,012           & 0,638 & 0,615 & 0,660 \\
\bottomrule
\end{tabular}
\end{table}

O Transformer exibe distribuição bimodal: os AUCs agrupam-se em dois extremos (sempre prever ``Sobe'' ou sempre prever ``Desce''), sem meio-termo discriminativo. O coeficiente de variação é de 38\%.

\section{Expanding-window cross-validation}

Substituindo o \textit{split} único por CV \textit{expanding-window} (janela mínima 600 dias, validação/teste 90 dias, passo 90 dias, 5 \textit{folds} $\times$ 5 sementes $\times$ 3 ativos = 145 treinamentos):

\begin{table}[htb]
\centering
\caption{AUC médio sob \textit{expanding-window} CV.}\label{tab:expandingcv}
\begin{tabular}{lccc}
\toprule
\textbf{Ativo} & \textbf{Baseline XGB} & \textbf{Transformer} & $\boldsymbol{\Delta}$ \\
\midrule
ITUB4 & \textbf{0,700} & 0,445 & $-$0,255 \\
PETR4 & \textbf{0,702} & 0,447 & $-$0,255 \\
VALE3 & 0,599           & 0,635 & +0,036 \\
\midrule
\textbf{Média} & \textbf{0,667} & 0,509 & $-$0,158 \\
\bottomrule
\end{tabular}
\end{table}

A inversão é quantitativamente expressiva ($d \approx 6{,}25$). O \textit{baseline} vence em 2 de 3 ativos por 0,25 pontos de AUC. A exceção aparente (VALE3, +0,036) é investigada em \textit{deep-dive} com 880 execuções e revelada como não significativa (Wilcoxon $p = 0{,}194$, IC~95\% $[-0{,}136;\ +0{,}033]$).

\section{Ablation: o sentimento adiciona algo?}

XGBoost em três configurações sob \textit{expanding-window} CV (3 ativos $\times$ 5 \textit{folds} $\times$ 5 sementes = 225 treinamentos):

\begin{table}[htb]
\centering
\caption{Ablation de \textit{features}: PRICE vs SENT vs PRICE+SENT.}\label{tab:ablation}
\begin{tabular}{lcccc}
\toprule
\textbf{Ativo} & \textbf{PRICE} & \textbf{SENT} & \textbf{PRICE+SENT} & $\boldsymbol{\Delta}$ \textbf{(P+S$-$P)} \\
\midrule
ITUB4 & 0,684 & 0,436 & 0,651 & $-$0,033 \\
PETR4 & 0,692 & 0,494 & 0,676 & $-$0,016 \\
VALE3 & 0,609 & 0,510 & 0,667 & +0,058 \\
\midrule
\textbf{Média} & \textbf{0,662} & 0,480 & 0,665 & \textbf{+0,003} \\
\bottomrule
\end{tabular}
\end{table}

O ganho incremental do sentimento é $\Delta = +0{,}003$ (Wilcoxon $p = 0{,}49$, IC~95\% $[-0{,}012;\ +0{,}018]$). A análise de poder estatístico (Hanley \& McNeil, 1982) confirma que este efeito está 44--74$\times$ abaixo do tamanho mínimo detectável para os tamanhos de teste utilizados (MDE entre 0,132 e 0,221 para $N$ entre 60 e 177).

\section{Validação do TCN Stage~5b}

O TCN~[32,32] com \textit{features} engenheiradas de sentimento --- melhor resultado prático (AUC = 0,643 sob janela única) --- foi submetido ao mesmo escrutínio: 20 sementes (média 0,513 $\pm$ 0,102) e \textit{expanding-window} CV (média 0,556). Sob nenhuma avaliação multi-fold o TCN supera o \textit{baseline} autoregressivo.

\section{Síntese}

A conclusão dos 13 experimentos é:

\begin{citacao}
Sob avaliação metodologicamente correta (multi-fold, multi-sementes, com intervalos de confiança e \textit{ablation} de \textit{features}), as \textit{features} de sentimento FinBERT-PT-BR utilizadas neste estudo não adicionam sinal preditivo mensurável a um \textit{baseline} autoregressivo de cinco \textit{features} de preço, em nenhum dos três ativos brasileiros de grande capitalização testados.
\end{citacao}


% ============================================================================
\chapter{Conclusão}\label{cap:conclusao}
% ============================================================================

\section{Contribuições}

Este trabalho oferece três contribuições:

\begin{enumerate}
    \item \textbf{Pipeline técnico replicável.} O sistema de coleta de notícias (InfoMoney), extração de sentimento (FinBERT-PT-BR) e integração com dados de mercado (Yahoo Finance) está integralmente versionado e pode ser reproduzido por outros pesquisadores.

    \item \textbf{Demonstração empírica do viés de janela única.} A inversão de AUC 0,709 (janela única) para $\sim$0,51 (multi-fold) é documentada com 1.435+ execuções e suporte estatístico formal. Esta é, no conhecimento do autor, a demonstração mais detalhada deste fenômeno em dados brasileiros.

    \item \textbf{Proposição de protocolos mínimos.} O trabalho propõe seis práticas mínimas para pesquisa em ML financeiro: (1)~reportar IC \textit{bootstrap}; (2)~treinar com $\geq$10 sementes; (3)~usar CV \textit{expanding-window}; (4)~comparar contra \textit{baseline} autoregressivo; (5)~monitorar distribuição de predições; (6)~auditar correlação validação-teste.
\end{enumerate}

\section{Limitações}

As principais limitações incluem: (1)~fonte única de notícias (InfoMoney), com viés editorial voltado ao investidor de varejo e cobertura heterogênea entre ativos; (2)~apenas 3 ações de grande capitalização da B3, limitando a validade externa; (3)~horizonte de previsão restrito a 5 e 21 dias úteis; (4)~a conclusão sobre o sentimento se aplica à representação específica adotada (5 \textit{features} agregadas diariamente do FinBERT-PT-BR), não a toda forma possível de codificar informação textual.

\section{Trabalhos futuros}

Cinco direções são sugeridas: (1)~diversificação de fontes textuais (comunicados CVM, \textit{analyst reports}, redes sociais); (2)~teste de arquiteturas menos propensas a colapso (modelos lineares com interações, GBT leve); (3)~generalização para outros mercados (S\&P~500, mercados emergentes); (4)~estudo formal da anticorrelação validação-teste em diferentes janelas; (5)~investigação de representações textuais alternativas (sentença-nível vs documento-nível, \textit{aspect-based sentiment}).


% ============================================================================
% ELEMENTOS PÓS-TEXTUAIS
% ============================================================================
\postextual

% --- Referências ---
\bibliography{referencias}

\end{document}
